\section{Stratégie de test et intégration continue}
\label{sec:tests}
\subsection{Des tests indépendants du service distant}
Les tests OpenAI utilisent le SDK réel face à un serveur HTTP local simulé ; ceux de Gemini injectent des fonctions de génération et de découverte contrôlées. Le test du moteur fournit une réponse LLM déterminée pour vérifier le contexte, la sélection du modèle et l'export. Ces tests ne nécessitent aucune clé API réelle.

Les tests de sécurité contrôlent les arguments Docker et l'exécution avec des doubles, ce qui les distingue des essais d'un vrai conteneur. Le moteur créant sa sandbox directement, son test peut tenter Docker, puis emprunter le repli statique si l'outil est absent. L'injection de cette dépendance rendrait ce scénario entièrement contrôlé.

\begin{table}[H]
\small
\begin{tabularx}{\textwidth}{@{}YP{19mm}P{64mm}@{}}
\toprule
\tabhead{Famille} & \tabhead{Réussis} & \tabhead{Comportements vérifiés}\\
\midrule
Configuration et objets métier & 4 & Catalogue, invariants et validité des données.\\
Moteurs et parseurs & 22 & Réponses attendues ou invalides, scénario d'analyse.\\
Fournisseurs LLM & 24 & Texte, erreurs HTTP, délai, modèle transmis, découverte.\\
Projet et contexte XML & 6 & Parcours, sélection des fichiers et limites.\\
Génération du rapport & 4 & Export, modèle documentaire et échappement.\\
Protection du prompt & 19 & Détection heuristique et délimitation du contenu.\\
Composants de sandbox & 25 & Configuration, restrictions, exécution simulée.\\
\midrule
\textbf{Total réussi} & \textbf{104} & \textbf{Aucun échec ni erreur dans la suite standard.}\\
\bottomrule
\end{tabularx}
\caption{Résultats de la dernière validation standard sur la version étudiée : 109 tests recensés, dont 5 tests JavaFX non exécutés. Les tests Docker nommés \code{*IT} ne sont pas inclus dans ce total.}
\end{table}

\subsection{Portée de la validation}
Les cinq tests graphiques requièrent une session graphique et une activation explicite. Les essais Docker se lancent séparément. Les compteurs du tableau proviennent des rapports Surefire ; ils ne constituent ni un taux de couverture du code ni une preuve de sécurité en conditions hostiles.

Une évaluation qualitative des appréciations du LLM sur des projets de référence reste nécessaire. Le choix des modèles, l'arrivée tardive de leur catalogue et les sorties de sandbox volumineuses méritent aussi une campagne dédiée.

\subsection{Intégration continue sur GitHub}
Le workflow \code{.github/workflows/test.yml} prévoit une validation sur les push et pull requests visant \code{main} ou \code{dev}, ainsi qu'un lancement manuel. Il prépare Java 21 et exécute \code{mvn verify}. Le cache Maven accélère les exécutions ; les droits sont limités à la lecture du dépôt. La configuration impose un délai maximal de quinze minutes et annule un ancien travail concurrent pour le même contexte de branche. Cette description porte sur le workflow versionné, sans préjuger du résultat d'une exécution distante particulière.

\clearpage
\section{Extensibilité et perspectives}
\label{sec:extensions}
\subsection{Points d'extension du prototype}
L'extensibilité est inégale selon les fonctionnalités. Les contrats couvrent déjà le fournisseur, le moteur et l'export ; la sélection des fichiers et la commande d'exécution restent plus directement codées dans leurs composants. Le tableau suivant précise le coût d'une extension sans supposer l'existence d'un mécanisme de plugins.

\begin{table}[H]
\small
\begin{tabularx}{\textwidth}{@{}P{32mm}YP{43mm}@{}}
\toprule
\tabhead{Évolution} & \tabhead{Point d'intervention} & \tabhead{Limite actuelle}\\
\midrule
Nouveau fournisseur & Implémenter \code{LlmProvider}, adapter requêtes et erreurs, puis compléter la fabrique. & La création reste centralisée dans un \code{switch}.\\
Nouveau critère & Ajouter une entrée à \code{criteria.json} et ses paramètres. & Ressource embarquée, rechargée au démarrage après mise à jour du build.\\
Nouveau format de rapport & Implémenter \code{ReportGenerator} et modifier l'assemblage ainsi que les libellés de l'UI. & Aucun choix de format dans la fenêtre actuelle.\\
Nouvelle règle de sélection & Modifier \code{XmlFileTreeBuilder}, ou extraire un contrat de politique de fichiers. & Règles codées en dur ; parcours visuel distinct.\\
Nouvelle méthode d'analyse & Implémenter \code{EvaluationEngine} avec le même résultat et le même observateur. & Pas de composition générale de plusieurs analyseurs.\\
Nouveau langage exécuté & Adapter l'image et la commande de sandbox au système de construction. & Le moteur lance actuellement une commande Maven.\\
\bottomrule
\end{tabularx}
\caption{Évolutions possibles et modifications nécessaires.}
\end{table}

\subsection{Priorités d'amélioration}
\textbf{Stabiliser les contrats.} La première priorité est d'aligner le prompt avec le schéma effectivement validé et de vérifier la compatibilité des modèles proposés avec le format de sortie demandé. La cohérence entre configuration affichée, paramètres de la tâche et rapport doit être testée lorsque la liste distante arrive en cours d'utilisation.

\textbf{Durcir les frontières d'entrée.} La collecte des journaux doit être bornée, filtrée et protégée contre les injections. Une préparation contrôlée des dépendances Maven et une distinction claire entre panne d'infrastructure et échec des tests rendraient les observations plus interprétables. L'environnement VM/Docker doit être préparé et testé explicitement.

\textbf{Étendre la couverture fonctionnelle.} L'import d'archives, un stockage persistant des analyses et l'accès à un LLM local répondraient à des besoins du sujet encore non couverts. La gestion de grands projets appelle ensuite un découpage du contexte et une consolidation des résultats intermédiaires. Ces mécanismes doivent préserver les preuves utilisées pour chaque appréciation.

\textbf{Mesurer la qualité des évaluations.} Une pondération effective, un jeu de projets étalons et des comparaisons répétées entre modèles permettraient d'étudier la stabilité des notes. La conformité JSON ne peut tenir lieu de validation scientifique des jugements produits.

\clearpage
\section{Bilan du projet}
\label{sec:bilan}
\subsection{Apports de la réalisation}
Le projet a transformé une intention d'\emph{analyse de code par IA} en responsabilités identifiables. L'interface, le moteur, le service externe et la production documentaire sont reliés par des contrats. Cette séparation facilite le développement parallèle et les tests sans dépendre d'une réponse réelle de modèle.

L'évolution des schémas a clarifié plusieurs ambiguïtés : le contrôleur graphique coordonne les actions, le moteur orchestre l'analyse, la sandbox exécute et l'export reçoit des données validées. Cette clarification permet de situer une erreur et d'identifier le composant chargé de la traiter.

Le sélecteur de modèles illustre ce découpage. La vue collecte une préférence, le contrôleur la copie dans une tâche, le moteur l'associe à l'analyse et l'adaptateur la transmet au SDK. Le rapport conserve le modèle demandé, sans que la vue connaisse la structure d'une requête Gemini.

\subsection{Compromis et enseignements}
La contrainte de temps a conduit à privilégier un parcours complet sur un répertoire local plutôt qu'une couverture exhaustive du sujet. Le prototype traite un contexte borné, conserve un résultat en mémoire et exporte du LaTeX. Il ne propose pas encore d'historique consultable, de fournisseur local livré, d'analyse distribuée ni de compilation PDF intégrée à l'interface.

Les design patterns ont surtout été utiles lorsqu'ils correspondaient à une variation réelle : adapter un SDK, substituer un moteur, recevoir une notification ou masquer la mécanique de Docker. À l'inverse, la configuration des critères montre qu'une modélisation par données peut être plus appropriée qu'une nouvelle hiérarchie de classes. La testabilité fournit ici un critère concret pour apprécier la pertinence d'une abstraction.

Une fenêtre réactive, une réponse valide, un export réussi et une appréciation pertinente sont quatre propriétés différentes. La dernière exige un protocole d'évaluation humain et expérimental qui dépasse les tests unitaires actuels.

\subsection{Usage des outils d'assistance}
Conformément au sujet, des outils d'assistance par IA ont été mobilisés pour accompagner l'intégration, l'analyse d'erreurs, la vérification d'API et la préparation documentaire. Ils ont également contribué à la rédaction et à la mise en forme du présent rapport à partir du brouillon du groupe et du code. La responsabilité des explications, des choix techniques et de la validation finale demeure collective. Aucune contribution individuelle supplémentaire n'est déduite du seul usage de ces outils.

\begin{pointcle}{Conclusion}
OctoGenere fournit une base fonctionnelle pour l'évaluation assistée de projets logiciels. Son apport essentiel est une architecture explicite autour d'un composant probabiliste, avec des points de substitution et de validation. La poursuite du travail doit désormais porter sur la cohérence des contrats, la maîtrise des sorties non fiables et la mesure de la qualité des évaluations.
\end{pointcle}

\clearpage
\section*{Annexe — Reproduction et repères techniques}
\addcontentsline{toc}{section}{Annexe — Reproduction et repères techniques}
\label{sec:annexe}
\subsection*{Lancer le prototype}
Les commandes suivantes sont exécutées à la racine du dépôt, avec un JDK 21 et Maven. Copier le modèle de configuration dans \code{.env}, puis y renseigner la clé et un identifiant de modèle valide pour le fournisseur choisi. Le fichier d'exemple partagé ne doit contenir aucune clé réelle.

\begin{lstlisting}
cp .env.example .env
mvn javafx:run
\end{lstlisting}

Dans la fenêtre, sélectionner le projet et les critères, vérifier le modèle retenu, puis lancer l'analyse. Une fois les résultats disponibles, le bouton d'export crée un dossier distinct sous \code{target/reports}. La documentation d'installation complémentaire figure dans le README du dépôt \cite{depot}.

\subsection*{Valider les composants}
\begin{lstlisting}
mvn --batch-mode --no-transfer-progress verify
\end{lstlisting}
Le résultat standard documenté est de 109 tests recensés : 104 réussis et 5 non exécutés. Les rapports XML sont produits dans \code{target/surefire-reports}. Le fichier \code{docs/rapport/validation-tests.json} en conserve une synthèse numérique pour la version décrite, sans les propriétés d'environnement des rapports bruts.

Les tests graphiques doivent être activés dans une session graphique :
\begin{lstlisting}
mvn -Doctogenere.ui.tests=true -Dtest=UiWorkflowTest test
\end{lstlisting}

\subsection*{Préparer les essais Docker}
Pour l'exécution de projets non fiables, préparer au préalable la machine virtuelle dédiée prescrite par le sujet, puis y construire l'image et lancer les essais d'intégration :
\begin{lstlisting}
docker build -t octogenere/sandbox:1.0 -f docker/Dockerfile docker
mvn -Dtest=fr.octogenere.security.sandbox.DockerSandboxExecutorIT test
\end{lstlisting}
Ces commandes sont un protocole de reproduction ; leur présence dans cette annexe ne signifie pas que l'environnement VM/Docker a été validé lors de la préparation du rapport.

\subsection*{Recompiler le présent rapport}
Avec XeLaTeX, latexmk et les polices TeX Gyre installés :
\begin{lstlisting}
latexmk -xelatex -interaction=nonstopmode -halt-on-error \
  -outdir=target/rapport Rapport-OctoGenere.tex
cp target/rapport/Rapport-OctoGenere.pdf Rapport-OctoGenere.pdf
\end{lstlisting}
Le document est distinct d'une évaluation automatiquement générée par l'application. Le logo institutionnel est inclus dans les ressources du rapport et les schémas sont construits en vectoriel avec TikZ ; aucune récupération réseau n'est nécessaire à la recompilation.

\clearpage
\renewcommand{\refname}{Références}
\addcontentsline{toc}{section}{Références}
\begin{thebibliography}{9}
\raggedright
\setlength{\itemsep}{11pt}
\bibitem{sujet}
Université Côte d'Azur, Master Informatique.
\emph{Project: AI Project Reviewer — Design of an Automated Evaluation Platform for AI-Assisted Software Projects}.
Sujet de projet, année universitaire 2026–2027.
Document fourni dans le dépôt : \code{docs/sujet2627-m1-en.pdf}. En particulier, sections 6, 8, 10, 11 et 13.

\bibitem{depot}
Groupe 3. \emph{OctoGenere : code source, documentation et tests}.
Branche \code{dev}, version \code{56e7ae8}, 11 septembre 2026.
\url{https://github.com/alexbertho/octogenere/tree/56e7ae8}.
Les descriptions d'implémentation et les résultats de test du rapport se rapportent à cette version.

\bibitem{gof}
Erich Gamma, Richard Helm, Ralph Johnson et John Vlissides.
\emph{Design Patterns: Elements of Reusable Object-Oriented Software}.
Addison-Wesley, 1994. ISBN 978-0-201-63361-0.
Référence terminologique pour les patrons de conception objets.

\bibitem{javafx}
OpenJFX. \emph{JavaFX 21 API — Task}.
Documentation officielle du modèle de tâches et des contraintes du thread JavaFX.
\url{https://openjfx.io/javadoc/21/javafx.graphics/javafx/concurrent/Task.html}.
Consultée le 11 septembre 2026.

\bibitem{gemini-models}
Google. \emph{Gemini API — Models}.
Documentation officielle des opérations \code{models.list} et \code{models.get} et des fonctionnalités annoncées par les modèles.
\url{https://ai.google.dev/api/models}.
Consultée le 11 septembre 2026.

\bibitem{docker}
Docker. \emph{Resource constraints}.
Documentation officielle des limites de ressources des conteneurs.
\url{https://docs.docker.com/engine/containers/resource_constraints/}.
Consultée le 11 septembre 2026.

\bibitem{ds4h}
Université Côte d'Azur, EUR DS4H.
\emph{Communication | Médias — Identité graphique de l'EUR DS4H}.
\url{https://ds4h.univ-cotedazur.fr/a-propos-de-leur-ds4h/communication-medias-1}.
Logo institutionnel téléchargé le 11 septembre 2026.
\end{thebibliography}

\vfill
\begin{pointcle}{Crédits iconographiques}
Le bandeau de couverture provient de la page de communication officielle de l'EUR DS4H \cite{ds4h}. Il est reproduit dans ses proportions et ses couleurs d'origine, avec les éléments institutionnels associés. Les schémas du corps du rapport ont été redessinés à partir des composants et interactions du code, en remplacement des représentations préliminaires du brouillon.
\end{pointcle}
